% !TEX root = ../main.tex
\section{Background and Related Work}
\label{sec:background}

\paragraph{Reasoning agents and speculative tool dispatch.}
A tool-using agent works in turns~\citep{yao2023react}.
In each turn the model reads the conversation so far, reasons, and emits a tool call.
The tool returns an observation, which is appended to the conversation and starts the next turn.
With a reasoning model, one turn can decode hundreds of tokens before the call appears.
The loop is serial: the tool cannot start before the call is written, and the next turn cannot start before the observation exists.

One line of work takes the tool off this critical path by predicting the call and executing it while the model is still generating~\citep{nichols2025speculative,sui2026paste,huang2025spagent,li2026specmemory,bai2026spork}.
In its self-speculative form, no extra model is needed~\citep{bai2026spork}.
A probe request continues the model's own reasoning prefix and is forced to write a tool call.
If the probe is confident, the call is executed at once.
The result is used only if the model later emits the same call, and the probe's tokens can serve as drafts when the model writes it.
We build on this mechanism and call it \emph{speculative tool dispatch}.
Its benefit is bounded by the tool latency, which is small when tools are fast and reasoning is long.

\paragraph{Cross-turn action speculation.}
A second line of work starts future steps before the current one is confirmed.
Interactive and Dynamic Speculative Planning let a cheaper approximation agent run ahead of the target agent, which verifies each step~\citep{hua2024isp,guan2025dsp}.
Speculative Actions formulates this propose--verify pattern for general agent actions and analyzes when it is lossless~\citep{ye2026speculative}.
LEAP trains a small drafter on the target's own actions, executes the drafted actions in a copy of the environment, and lets the target verify the resulting contexts concurrently~\citep{xu2026leap}.
Other systems speculate multi-call macros~\citep{liu2026smc}, combine fast and slow reasoning modes~\citep{zhong2026dualspec}, or predict the observation instead of waiting for the tool~\citep{chen2026aospec,montgomery2026obsspec}.

These methods obtain the future input from a separate component: a cheaper or trained proposer for the action, or a learned model for the observation.
Their gain depends on how often this component agrees with the target, and a trained proposer must be fitted to each workload.
\method needs no such component.
It takes the action from the target's own unfinished turn and the observation from actually executing that action.
It also commits a turn only when its entire input is unchanged, which targets identity of whole trajectories rather than of single actions (\cref{sec:method:commit}).

\paragraph{Faster reasoning.}
Speculative decoding accelerates generation inside a turn at the token level and preserves the model's output~\citep{leviathan2023fast,chen2023accelerating}.
It is orthogonal to \method.
ETA observes, as we do in \cref{sec:motivation}, that reasoning continues long after the dispatched tool has returned, and it ends the reasoning early once the observation arrives~\citep{anon2026eta}.
This shortens the turn but changes what the model generates.
\method keeps the reasoning intact and uses the remaining time to run the next turn.
